% 04 - MACHINE D'ETATS
% ==============================================================
\sectionintro{04 · Séquencement}{Machine d'états du cycle AUTO}{La représentation reprend les états et les transitions réellement présents dans le bloc fonctionnel.}

\twocolpage{
  \begin{center}
  \begin{tikzpicture}[
    node distance=8mm,
    state/.style={draw=ekprimarydark,very thick,fill=white,rounded corners=1mm,
      minimum width=38mm,minimum height=8mm,font=\bfseries\scriptsize,
      text=ekprimarydark,align=center}
  ]
    \node[state,double,double distance=1.2pt,fill=eksoft] (a0) {A0 · ATTENTE};
    \node[state,below=of a0,fill=eksoft] (a1) {A1 · COMPTAGE};
    \node[state,below=of a1] (a2) {A2 · BAC\_PLEIN};
    \node[state,below=of a2] (a3) {A3 · VERIN\_SORTIE};
    \node[state,below=of a3] (a4) {A4 · VERIN\_RENTREE};
    \node[state,below=of a4] (a5) {A5 · BAC\_SUIVANT};
    \node[state,below=of a5,fill=ekwarning] (a6) {A6 · ATTENTE\_BAC};
    \draw[flow] (a0)--(a1);
    \draw[flow] (a1)--(a2);
    \draw[flow] (a2)--(a3);
    \draw[flow] (a3)--(a4);
    \draw[flow] (a4)--(a5);
    \draw[flow] (a5)--node[right,font=\scriptsize,text=ekmuted] {bac plein}(a6);
    \draw[flow] (a4.west)--++(-8mm,0)|-(a3.west);
    \draw[flow] (a5.west)--++(-13mm,0)|-(a1.west);
    \draw[flow] (a6.east)--++(7mm,0)|-(a1.east);
  \end{tikzpicture}
  \end{center}
}{
  \subsection{Transitions principales}
  \begin{tabularx}{\linewidth}{@{}L{24mm}Y@{}}
    \toprule
    \textbf{Transition} & \textbf{Condition}\\
    \midrule
    A0 → A1 & premier traitement de \texttt{ATTENTE} en mode \texttt{AUTO} ; transition sans condition\\
    A1 → A2 & \texttt{aPieceBac[nBac] >= nPieceBac}\\
    A2 → A3 & transition immédiate\\
    A3 → A4 & temporisation de sortie écoulée\\
    A4 → A3 & temporisation de rentrée écoulée et \texttt{nImp < 4}\\
    A4 → A5 & temporisation écoulée et \texttt{nImp >= 4}\\
    A5 → A1 & nouveau bac sous la consigne\\
    A5 → A6 & nouveau bac déjà plein\\
    A6 → A1 & \texttt{aResetBac[nBac] = TRUE}\\
    \bottomrule
  \end{tabularx}

  \begin{primarycard}
    \textbf{Lecture du cycle vérin}\par\small
    Une unité de \texttt{nImp} correspond à une phase avant temporisée suivie d'une phase arrière temporisée. Le numéro logique de bac ne change qu'après quatre cycles complets ; le code ne vérifie pas la position mécanique du plateau.
  \end{primarycard}

  \begin{notebox}[colback=ekwarning,colframe=orange!35]{Référencement}
    \texttt{INIT} exploite uniquement l'entrée logique \texttt{bZero}. La nature physique du capteur et la stratégie mécanique de recherche du zéro ne sont pas définies dans les sources ; \texttt{bZero = TRUE} valide le référencement.
  \end{notebox}
}

\newpage

% ==============================================================
